\section{Method: the \R Skills}
\label{sec:system}

\R is four agent skills and an interface. A skill packages instructions for the agent, the programs
that do the work, and acceptance checks that decide whether a step succeeded. The agent chooses which
step to run and with which arguments; it never certifies a result. Every handoff in
\cref{fig:overview} is therefore gated by a program, and every run, failed or not, leaves an evidence
record.

\subsection{Environment provisioning (eda-install)}
\label{sec:skill-install}
The skill turns an empty machine into a verified toolchain in four steps: \emph{detect} what the host
already provides (operating system, C library, Python, existing tools and kits), \emph{plan} the
missing components as tiers, \emph{install} them, and \emph{pin} the resolved paths in a single
environment file that every other skill reads. The tiers are: the place-and-route and synthesis
engines, distributed as a relocatable bundle pinned by checksum to the exact build used for all our
data; the open-source RTL-to-GDS flow at the commit validated with that bundle; simulators, layout
extraction and LVS tools; the process design kits and their signoff rule decks; and a graph-learning
Python environment. Installation needs no administrator rights, and a dry run prints the full plan
before anything is changed. Verification is per component: each tool must report the expected
version and pass a smoke check, and each kit counts as ready only when a reference design completes
the full flow with clean DRC and LVS on it.

\subsection{RTL acquisition (rtl-acquire)}
\label{sec:skill-rtl}
The skill supplies designs at corpus scale. \emph{Discovery} searches public code hosts with frozen
seed queries, or reads local trees and repository manifests. \emph{Screening} infers each candidate's
top module and its compile closure (every included file, package and macro), and rejects candidates
whose closure is incomplete or whose provenance is unclear. \emph{Qualification} runs synthesis only;
a design is qualified if it synthesises to a non-trivial netlist. Qualified designs are deduplicated
by structural hash against everything seen before, scored for quality, and exported as pre-layout
netlist graphs. Every rejection is recorded with its reason, which gives the per-handoff yield, and
synthesis failures are recorded as learnable events. Finally, a qualified candidate is promoted to a
full physical-design project in one step.

\subsection{Physical execution and repair (signoff-loop)}
\label{sec:skill-signoff}
The skill drives each project through synthesis, floorplanning, placement, clock-tree synthesis,
routing and full-deck signoff: design-rule checking, layout-versus-schematic comparison and parasitic
extraction with each kit's official decks. A sample counts as signed off only if DRC is clean and LVS
matches on the same final layout.

\paragraph{Repair.}
When a stage or a check fails, the failure is reduced to a \emph{symptom signature}: the failing
check, the violation class and a set of predicates over the reports. Signatures, not design names,
index a store of repair recipes, so a recipe learned on one design or kit can apply to another with
the same symptom. A recipe may change physical-implementation parameters only; the RTL, clock period,
timing constraints, process kit and signoff decks are fixed.

\paragraph{Learning and freezing.}
Every run, including failed and abandoned repair attempts, is ingested into a knowledge store of
outcomes, kept separate from a journal of the actions taken. Candidate recipes mined from the store
enter a lifecycle (shadow $\rightarrow$ candidate $\rightarrow$ promoted, or demoted) and are
promoted only after A/B trials against a control arm, judged with a variance-aware lower confidence
bound over repeats. For deployment the promoted set is \emph{frozen}: the knowledge is read-only, so
results are reproducible and the cost of covered failures falls to recipe application.

\subsection{Graph dataset construction (def-graph)}
\label{sec:skill-graph}
The skill converts a signed-off layout into training data, and it refuses layouts whose signoff
evidence is missing or does not belong to that run. From the final layout, parasitics, cell libraries
and technology files it extracts per-cell and per-net \emph{labels} (routing congestion, wire length,
path timing, IR drop, and parasitic resistance and capacitance) and \emph{features} for gates, nets,
pins and I/O ports. It then emits five graph views, from the bipartite gate--pin--net graph to a
gate-only graph with nets folded into edges, as heterogeneous graphs.

\paragraph{Stage-aware datasets.}
For prediction at the floorplan, placement, clock-tree and routing stages, the skill builds one graph
per stage over a shared node set taken from the synthesized netlist, so circuit identity is stable
across stages. Each stage's features are restricted to information available at that stage
(\emph{feature visibility}), while labels come from the routed result with explicit validity masks.

\paragraph{Independent verification.}
A separate verifier recomputes every label and feature from the raw layout, libraries and checker
reports with code independent of the extractor. It checks graph topology, feature statistics and
label agreement with the signoff reports, and rejects any stage graph that uses information from a
later stage.

\subsection{Interface}
\label{sec:skill-gui}
A web interface exposes the workflow to users who do not read logs. A user submits a design,
chooses the platform and the language model (any provider with a compatible API, configured by the
user), and follows the run on a live timeline of stages and repair attempts. Each stage view shows
its layout image, and the run view shows the signoff verdicts, the model calls and tokens used, and
the recipes applied. Any view can be exported as a single-page PDF figure (\cref{fig:ui}).

\begin{figure}[t]
  \centering
  \PH{Interface screenshot, exported as a single-page PDF}
  \caption{The \R interface during a repair run: stage timeline, layout view and signoff verdicts.
  \PH{final caption}}
  \label{fig:ui}
\end{figure}
